%% kapa: exam=2 round=23 year=2012 period=3 subject=compensation_law subject_ko=감정평가 및 보상법규 seq=1 number=1 points=40
%% source: Q-Net 감정평가사 기출문제 (cst00309) artlSeq=5056014 file=raw/1008/5056014/3교시(감정평가 및 보상법규).hwp
%% conversion: kordoc 4.24.51 -> tools/gen2.py (auto); review=proofread-sonnet; flags=none
\begin{longquestion}{【문제 1】}
% [원문 그대로] "2011월 4월" — 시험지 원문 표기(오탈자 추정)를 그대로 둠
A도는 2008년 5월경 국토해양부장관으로부터 관계 법령에 따라 甲의 농지 4,000m\textsuperscript{2}를 포함한 B시와 C시에 걸쳐있는 토지 131,000m\textsuperscript{2}에 ‘2009 세계엑스포’ 행사를 위한 문화시설을 설치할 수 있도록 하는 공공시설입지승인을 받았다. 그 후 A도는 편입토지의 소유자들에게 보상협의를 요청하여 甲으로부터 2008년 12월 5일 「공익사업을 위한 토지 등의 취득 및 보상에 관한 법률」에 의하여 위 甲의 농지를 협의취득하였다. A도는 취득한 甲의 토지 중 1,600m\textsuperscript{2}를 2009년 5월 31일부터 2011월 4월 30일까지 위 세계엑스포행사 및 기타 행사를 위한 임시주차장으로 이용하다가 2012년 3월 31일 농지로 원상복구하였다. 그 후 위 1,600m\textsuperscript{2}의 토지는 인근에서 청소년수련원을 운영하는 제3자에게 임대되어 청소년들을 위한 영농체험 경작지로 이용되고 있다. (40점) \par
1) 甲은 농지로 원상복구된 토지 1,600m\textsuperscript{2}에 대한 환매권을 행사하려고 한다. 甲의 권리구제방법에 대하여 설명하시오. (25점) \par
2) A도는 환매권 행사 대상 토지의 가격이 현저히 상승된 것을 이유로 증액된 환매대금과 보상금상당액의 차액을 선이행하거나 동시이행할 것을 주장하려 한다. 환매대금 증액을 이유로 한 A도의 대응수단에 대하여 설명하시오. (15점) \par
\end{longquestion}
